\begin{table}[!htbp]
\centering
\caption{Continuous-bullying robustness model}
\label{tab:bullying-robustness}
\scriptsize
\setlength{\tabcolsep}{2.2pt}
\resizebox{\textwidth}{!}{%
\begin{tabular}{lcc}
\toprule
 & \multicolumn{1}{c}{South Africa} & \multicolumn{1}{c}{Brazil} \\
\cmidrule(lr){2-2}\cmidrule(lr){3-3}
 & R1 & R1 \\
\midrule
Bullying scale (higher = less bullying) & \shortstack{18.7***\\(1.3)} & \shortstack{16.8***\\(1.2)} \\
Boy & \shortstack{-43.7***\\(3.6)} & \shortstack{-5.0\\(4.8)} \\
Age (years) & \shortstack{-8.0**\\(3.4)} & \shortstack{-13.7***\\(4.9)} \\
Home SES (ASBHSES scale points) & \shortstack{26.1***\\(2.3)} & \shortstack{27.4***\\(2.1)} \\
Other gender category &  & \shortstack{-97.7**\\(36.8)} \\
\midrule
Learners & 9,167 & 4,110 \\
$R^{2}$ & 0.270 & 0.293 \\
\bottomrule
\end{tabular}%
}
\parbox{0.98\textwidth}{\footnotesize \textit{Notes:} Entries are coefficients with jackknife standard errors in parentheses. Higher ASBGSB values indicate less bullying. Significance: *** p\textless{}0.01, ** p\textless{}0.05, * p\textless{}0.10.}
\end{table}
